% Part: normal-modal-logic
% Chapter: axioms-systems
% Section: introduction

\documentclass[../../../include/open-logic-section]{subfiles}

\begin{document}

\olfileid{nml}{axs}{int}

\olsection{పరిచయం}

ప్రాథమిక మోడల్ భాషకు మోడల్ నమూనాల ద్వారా అర్థవిచారం మనకు ఉంది.
\tetoken{ఒక సూత్రం}{!!a{formula}} చెల్లుబాటవుతుందనే భావన కూడా ఉంది:
అది అన్ని నమూనాల్లోని అన్ని లోకాల వద్ద సత్యం. అలాగే ఏదైనా నమూనాల
లేదా చట్రాల వర్గానికి సంబంధించి చెల్లుబాటవుతుందనే భావన ఉంది:
ఆ వర్గంలోని అన్ని నమూనాల్లోని అన్ని లోకాల వద్ద, లేదా ఆ చట్రంపై
ఆధారపడిన అన్ని నమూనాల్లోని అన్ని లోకాల వద్ద అది సత్యం.
తర్కం సాధారణంగా చెల్లుబాటు గురించిన ఇటువంటి అర్థవిచార వర్ణనలను
\tetoken{వ్యుత్పాద్యత}{!!{derivability}} అనే నిరూపణ-సిద్ధాంత భావనతో
అనుసంధానిస్తుంది. \tetoken{ఒక సూత్రం}{!!a{formula}} చెల్లుబాటైతే,
అప్పుడే, అప్పుడు మాత్రమే, అది \tetoken{వ్యుత్పాదించదగినది}{!!{derivable}}
అయ్యేలా ఏదో వ్యవస్థలో \tetoken{వ్యుత్పాద్యత}{!!{derivability}}ను
నిర్వచించడమే లక్ష్యం.

అత్యంత సరళమైన, చారిత్రకంగా అతి పాత \tetoken{వ్యుత్పత్తి}{!!{derivation}}
వ్యవస్థలు హిల్బర్ట్-రకం లేదా స్వీకృతాధారిత \tetoken{వ్యుత్పత్తి}{!!{derivation}}
వ్యవస్థలు. అనేక మోడల్ తర్కాల కోసం హిల్బర్ట్-రకం
\tetoken{వ్యుత్పత్తి}{!!{derivation}} వ్యవస్థలను నిర్మించడం సాపేక్షంగా
సులభం: అధిసిద్ధాంత పరిశీలనకు అవి సరళమైన వస్తువులు (ఉదాహరణకు,
నిర్దుష్టతను, సంపూర్ణతను నిరూపించడానికి). అయితే, సహజ నిగమన
వ్యవస్థలతో పోలిస్తే వాటిలో \tetoken{సూత్రాలను}{!!{formula}s}
నిరూపించడం చాలా కష్టం.

హిల్బర్ట్-రకం \tetoken{వ్యుత్పత్తి}{!!{derivation}} వ్యవస్థల్లో,
\tetoken{ఒక సూత్రం}{!!a{formula}} యొక్క \tetoken{వ్యుత్పత్తి}{!!a{derivation}}
అంటే కొన్ని స్వీకృతాల నుంచి కొద్దిపాటి నిగమన నియమాల ద్వారా
ప్రశ్నలోని \tetoken{సూత్రానికి}{!!{formula}} చేరే \tetoken{సూత్రాల}{!!{formula}s}
క్రమం. \tetoken{వ్యుత్పత్తి}{!!{derivation}} వ్యవస్థ అర్థవిచారానికి
సరిపోవాలని కోరుతున్నాం కనుక, \tetoken{వ్యుత్పాదించదగిన}{!!{derivable}}
సూత్రాల సమితి అన్ని నమూనాల్లో సత్యమని (లేదా అన్ని స్వీకృతాలూ
సత్యమయ్యే ప్రతి నమూనాలో సత్యమని) నిర్ధారించాలి. ఇందుకు
అవసరమైన మోడల్ తర్కాల కొన్ని ధర్మాలను ముందుగా వేరుచేస్తాం:
అవే ``నార్మల్'' మోడల్ తర్కాలు. నార్మల్ మోడల్ తర్కాల కోసం
మోడస్ పోనెన్స్, అవశ్యకీకరణ అనే రెండు నిగమన నియమాలను మాత్రమే
అనుమానించాలి. స్వీకృతాలుగా సర్వసత్యాల అన్ని ప్రతిస్థాపన
నిదర్శనాలను, మనం పరిశీలించే మోడల్ తర్కాన్ని బట్టి కొన్ని మోడల్
స్వీకృతాలను తీసుకుంటాం. అన్ని నమూనాల వర్గంపైనే మనకు ఆసక్తి
ఉన్నప్పటికీ,~$\Ax{K}$\iftag{notprvDiamond}{}{ మరియు~$\Ax{Dual}$} యొక్క
అన్ని ప్రతిస్థాపన నిదర్శనాలను కూడా స్వీకృతాలుగా లెక్కించాలి.
ఈ నియమాలు, స్వీకృతాల ఏర్పాటే కనిష్ఠ నార్మల్ మోడల్
తర్కం~$\Log K$ను ఉత్పత్తి చేస్తుంది.

\begin{defn}
\emph{మోడస్ పోనెన్స్} నియమం కింది నిగమన పథకం:
\begin{prooftree}
\AxiomC{$!A$}
\AxiomC{$!A \lif !B$}
\RightLabel{\MP}
\BinaryInfC{$!B$}
\end{prooftree}
\tetoken{ఒక సూత్రం}{!!a{formula}}~$!B$, $!A$, $!C$ అనే
\tetoken{సూత్రాల}{!!{formula}s} నుంచి మోడస్ పోనెన్స్ ద్వారా
\emph{అనుసరిస్తుంది} అని అంటాం, అప్పుడే, అప్పుడు మాత్రమే,
$!C \ident !A \lif !B$.
\end{defn}

\begin{defn}
\emph{అవశ్యకీకరణ} నియమం కింది నిగమన పథకం:
\begin{prooftree}
\AxiomC{$!A$}
\RightLabel{\Nec}
\UnaryInfC{$\Box !A$}
\end{prooftree}
\tetoken{సూత్రం}{!!{formula}}~$!B$, \tetoken{సూత్రాల}{!!{formula}s}లోని $!A$
నుంచి అవశ్యకీకరణ ద్వారా అనుసరిస్తుంది అని అంటాం, అప్పుడే,
అప్పుడు మాత్రమే, $!B \ident \Box !A$.
\end{defn}

\begin{defn}
స్వీకృతాల సమితి~$\Sigma$ నుంచి \emph{\tetoken{వ్యుత్పత్తి}{!!{derivation}}}
అంటే $!B_1$, $!B_2$, \dots, $!B_n$ అనే \tetoken{సూత్రాల}{!!{formula}s}
క్రమం; ఇందులో ప్రతి $!B_i$ కింది వాటిలో ఏదో ఒకటి:
\begin{enumerate}
\item ఒక సర్వసత్యానికి ప్రతిస్థాపన నిదర్శనం; లేదా
\item $\Sigma$లోని \tetoken{ఒక సూత్రానికి}{!!a{formula}} ప్రతిస్థాపన
  నిదర్శనం; లేదా
\item $j$, $k < i$ అయ్యే $!B_j$, $!B_k$ అనే రెండు
  \tetoken{సూత్రాల}{!!{formula}s} నుంచి మోడస్ పోనెన్స్ ద్వారా అనుసరించేది; లేదా
\item $j < i$ అయ్యే \tetoken{ఒక సూత్రం}{!!a{formula}}~$!B_j$ నుంచి
  అవశ్యకీకరణ ద్వారా అనుసరించేది.
\end{enumerate}
$!B_n \ident !A$ అయ్యే ఇటువంటి \tetoken{ఒక వ్యుత్పత్తి}{!!a{derivation}}
ఉంటే, $!A$ను \emph{$\Sigma$ నుంచి \tetoken{వ్యుత్పాదించదగినది}{!!{derivable}}}
అంటాం; సంకేతంగా $\Sigma \Proves !A$.
\end{defn}

ఈ నిర్వచనం ప్రకారం \tetoken{వ్యుత్పాదించదగిన}{!!{derivable}} సూత్రాల
సమితి ఒక నార్మల్ మోడల్ తర్కాన్ని ఏర్పరుస్తుందని, అన్ని స్వీకృతాలూ
సత్యమయ్యే ప్రతి నమూనాలో ప్రతి \tetoken{వ్యుత్పాదించదగిన}{!!{derivable}}
\tetoken{సూత్రం}{!!{formula}} సత్యమని తేలుతుంది. \tetoken{వ్యుత్పత్తుల}{!!{derivation}s}
ఈ ధర్మాన్ని \emph{నిర్దుష్టత} అంటారు. దీని విలోమమైన
\emph{సంపూర్ణత}ను నిరూపించడం మరింత కష్టం.

\end{document}
